\section{WireGuard}\label{wireguard}

La \textbf{VPN} è un servizio usato per comunicare in modo sicuro tra
membri di un gruppo attraverso l'uso di una infrastruttura di
comunicazione pubblica su cui il traffico viene incapsulato dentro ad un
tunnel cifrato. \footnote{\label{fn:ref_27}M. S. et al., \emph{Virtual Private Network
  (VPN) Protocols and Architecture}, IEEE,
  https://ieeexplore.ieee.org/document/4756450}

\textbf{Wireguard} è una VPN che utilizza il concetto di
\textbf{cryptokey routing}: viene mantenuta in memoria una tabella dove
ad ogni chiave pubblica di un peer, viene associata una lista di
indirizzi ip identificati come allowed.\footnote{\label{fn:ref_28}J. A. Donenfeld,
  \emph{WireGuard: Next Generation Kernel Network Tunnel}, NDSS
  Symposium, 2017. https://www.wireguard.com/papers/wireguard.pdf}
Fornisce un'interfaccia stateless per utilizzare il tunnel VPN, ossia
un'interfaccia di rete virtuale, che potrebbe essere per esempio
\passthrough{\lstinline!wg0!}.

\begin{quote}
Per \textbf{stateless}, si intende il fatto che l'amministratore non
deve occuparsi della gestione esplicita dello stato di sessione,
dell'handshake periodico o della riconnessione, poiché tali compiti sono
demandati automaticamente al modulo kernel di WireGuard.
\end{quote}

Rispetto ad OpenVPN, WireGuard è considerato superiore sul piano della
sicurezza e dell'efficienza:

\begin{itemize}
\tightlist
\item
  OpenVPN supporta una vasta gamma di cifrari e opzioni TLS che
  incrementano notevolmente la superficie di attacco complessiva.
\item
  OpenVPN opera nello spazio utente tramite dispositivi virtuali
  TUN/TAP, introducendo un overhead significativo dovuto a molteplici
  passaggi di contesto e copie di buffer di memoria tra \emph{user
  space} e \emph{kernel space}.
\end{itemize}

WireGuard non implementa un protocollo applicativo per lo scambio e la
negoziazione iniziale delle chiavi, delegando tale incombenza a canali
di distribuzione sicuri esterni e preesistenti (quali ad esempio
orchestratori, agenti software dedicati o canali SSH protetti).

\pagebreak

La struttura concettuale di una \emph{cryptokey routing table} associa a
ciascun peer autorizzato la relativa chiave pubblica crittografica e la
lista di indirizzi IP sorgente consentiti:

\begin{longtable}[]{@{}
  >{\raggedright\arraybackslash}p{(\linewidth - 6\tabcolsep) * \real{0.2500}}
  >{\raggedright\arraybackslash}p{(\linewidth - 6\tabcolsep) * \real{0.2500}}
  >{\raggedright\arraybackslash}p{(\linewidth - 6\tabcolsep) * \real{0.2500}}
  >{\raggedright\arraybackslash}p{(\linewidth - 6\tabcolsep) * \real{0.2500}}@{}}
\toprule\noalign{}
\begin{minipage}[b]{\linewidth}\raggedright
Peer Public Key
\end{minipage} & \begin{minipage}[b]{\linewidth}\raggedright
Allowed Source IPs
\end{minipage} & \begin{minipage}[b]{\linewidth}\raggedright
Endpoint (Opzionale)
\end{minipage} & \begin{minipage}[b]{\linewidth}\raggedright
Descrizione
\end{minipage} \\
\midrule\noalign{}
\endhead
\bottomrule\noalign{}
\endlastfoot
\passthrough{\lstinline!xTIB...p8Dg!} &
\passthrough{\lstinline!10.k.0.2/32!} &
\passthrough{\lstinline!198.51.100.1:51820!} & Giocatore 1 (Team
\(k\)) \\
\passthrough{\lstinline!TrMv...WXX0!} &
\passthrough{\lstinline!10.k.0.3/32!} &
\passthrough{\lstinline!203.0.113.2:51820!} & Giocatore 2 (Team
\(k\)) \\
\passthrough{\lstinline!gN65...z6EA!} &
\passthrough{\lstinline!10.k.0.4/32!} &
\passthrough{\lstinline!198.51.100.5:51820!} & Giocatore 3 (Team
\(k\)) \\
\end{longtable}

La tabella determina quali coppie
\((\text{Peer Public Key, Allowed Source IP})\) abbiano facoltà di
inviare e ricevere pacchetti attraverso l'interfaccia virtuale
\passthrough{\lstinline!wg-k!}.

Opzionalmente è possibile specificare verso quale endpoint inoltrare il
pacchetto. L'\textbf{endpoint} identifica l'indirizzo IP pubblico e la
porta UDP su cui il demone o modulo WireGuard remoto è in ascolto.

Il flusso di elaborazione e instradamento dei pacchetti si articola
nelle seguenti fasi:

\begin{itemize}
\tightlist
\item
  l'interfaccia \passthrough{\lstinline!wg0!} riceve un pacchetto IP da
  inoltrare nel tunnel;
\item
  il kernel incapsula il pacchetto all'interno di un datagramma UDP il
  cui payload viene autenticato e cifrato mediante il cifrario
  simmetrico \passthrough{\lstinline!ChaCha20Poly1305!};
\item
  il pacchetto cifrato viene trasmesso verso l'endpoint specificato,
  oppure verso l'indirizzo remoto appreso dinamicamente in occasione
  dell'ultimo handshake;
\item
  il nodo ricevente valida la corrispondenza dell'indirizzo sorgente
  nella \emph{cryptokey routing table};
\item
  il payload viene decifrato e autenticato, estraendo il pacchetto
  originale;
\item
  il pacchetto in chiaro viene immesso nello stack di rete locale
  attraverso l'interfaccia virtuale \(\text{wg}0\).
\end{itemize}
